\chapter{The Security and Governance Layer}

Agents introduce a category of risk that conversational systems do not. A model that
produces an incorrect response misleads a user. An agent that issues an incorrect
instruction to a tool can leak sensitive data, send communications, modify records, or
deploy broken code to production. Security and governance are therefore foundational
properties of the architecture, rather than capabilities added after the core system is
built.

\section{The Attack Surface}

The concept of an attack surface describes the set of entry points through which an
adversary can influence a system's behaviour. In agentic systems, this surface is
substantially wider than in conventional applications. The model's context window is an
attack vector: every piece of content that enters it---retrieved documents, tool outputs,
emails, web pages, tickets, and code---can contain instructions that attempt to redirect
the model's behaviour. The agent's memory is a further attack vector: content written
into long-term memory during one session can influence the model's reasoning in future
sessions, a class of attack known as memory poisoning. The tools themselves constitute a
third vector: an agent with access to credentials, external APIs, or write operations
over persistent storage can be induced to exercise those capabilities in ways that were
not authorised.

The principle of least privilege is the foundational response to this expanded attack
surface. Each component of the system---the agent, each tool, each memory store---should
hold only the permissions required to perform its declared function, scoped to the
current task, and enforced by the infrastructure rather than assumed from the model's
compliance with its system prompt. In agentic contexts, least privilege means limiting
an agent's access to only the resources and actions explicitly required to complete its
current task, not its general purpose, its potential future tasks, or everything it might
conceivably need.

\section{Prompt Injection}

Prompt injection occurs when content from an untrusted source attempts to redirect the
model's behaviour. Direct injection arrives in the user's own message. Indirect injection
is embedded in content the agent retrieves or processes---a web page, a PDF, an email
body, a repository file, or a tool output---and is particularly consequential because the
user may have no visibility into the injected instruction.

The primary structural mitigations are architectural. Context should be labelled by
source and authority level so the model can distinguish between a user instruction and a
retrieved passage. Tool access should be scoped to the authenticated user's identity so
that an injected instruction cannot invoke a capability the user does not hold. Tool
calls should be validated by an independent component rather than trusted on the basis
of the model's output alone. Data flows from sensitive sources to external sinks should
be explicitly governed. Write operations should require approval for consequential
actions. High-risk tools should operate against an allowlist. Every consequential action
should be logged with sufficient context to reconstruct the decision that produced it.
Infrastructure enforcement is what gives these controls their guarantee.

\section{Data Flow Control}

A governed agent architecture tracks the provenance of information from source to sink.
A source is where information enters the agent's context: a user profile, a private
document, an email, a code repository, a database record, or a web page. A sink is where
information leaves the system: a final response, an outbound email, a message, a ticket
update, a file write, or an external API call.

Data exfiltration occurs when content from a protected source reaches a sink the user or
organisation has not authorised. An agent with read access to internal documents and
write access to external communication channels must be governed by explicit data flow
rules that an injected instruction cannot override. These rules belong in the
infrastructure layer, enforced at the tool boundary.

\section{Human Approval}

For actions that are irreversible or consequential---modifying or deleting records,
deploying code, submitting a financial or legal document, or sharing internal data
outside the organisation---a human approval step is one available control. The
appropriate level of human oversight depends on the risk profile of the action and the
degree of automation the use case requires. A fully automated pipeline may accept higher
risk in exchange for throughput; a governed enterprise deployment may route every write
action through an approval queue. The design choice is between automation and
accountability, and the right point on that spectrum is determined by the consequences
of an unchecked error.

Approval events carry evaluation value regardless of where that point is set. A rejected
approval records a case where the agent's proposed action diverged from human judgment.
Collected systematically, rejected approvals become labelled examples for the evaluation
dataset and regression tests for subsequent agent versions.

\section{Sandboxing}

A sandbox is an environment that replicates production conditions without connecting to
production systems. In a staged deployment pipeline, the sandbox stage---typically
referred to as gamma---is where the agent is tested against realistic data and realistic
tool schemas before any change is promoted to production. The CI/CD pipeline should
include a sandbox stack where prompt changes, tool schema updates, and new agent
capabilities can be exercised and evaluated in isolation. Capabilities validated in the
sandbox environment carry a reproducible evidence record; those that have not been
validated there should remain outside the production deployment.

\section{Governance and Auditability}

Governance requires that the organisation can reconstruct, for any past interaction, who
invoked the agent, what the model decided, which tools were called, and which policy
checks ran. Observability infrastructure that captures every model call, tool invocation,
and policy check as a named span with associated metadata is the prerequisite for this.
Agent versioning ensures that a trace can be associated with the exact prompt, tool
schemas, and skill documents that were active at the time of the interaction, making it
possible to reproduce and evaluate historical behaviour as the system evolves.

\section{Summary}

The security and governance layer enforces the boundaries within which the agent
operates. The attack surface of an agentic system spans the context window, the memory
layer, and the tool interfaces; the principle of least privilege is the foundational
response to each. Prompt injection, data flow control, human approval calibrated to the
risk profile of each action, sandboxed validation before promotion, and structured
auditability are the core concerns of this layer. Each of them is an infrastructure
property enforced independently of the model, and each connects directly to the
evaluation and deployment practices covered in subsequent chapters.